\bookchapter{제52장: 안드로이드 4대 컴포넌트와 Activity 생명주기}

\chapterintro{Compose가 앱의 “얼굴”이라면, 이번 장은 앱의 “뼈대” 이야기예요. 안드로이드 시스템에는 4대 기본 컴포넌트가 있고 모든 앱은 이것들로 조립돼요. 그중 가장 중요한 Activity에는 모든 안드로이드 개발자가 반드시 머리에 새겨 넣어야 할 “생명주기 그림”이 있죠. 이 장을 끝내면 “뒤로 나갔다 돌아왔을 뿐인데 왜 내 앱 상태가 사라졌을까?” 같은 오랜 난제를 이해하게 될 거예요.}

\booksection{4대 컴포넌트: 안드로이드의 “기둥과 보”}

안드로이드는 각각 한 가지 일을 맡는 네 가지 기본 컴포넌트를 공식적으로 정의하고 있어요:

\begin{center}
\begin{tabularx}{\linewidth}{|X|X|X|}
\hline
\textbf{컴포넌트} & \textbf{역할} & \textbf{비유하면} \tabularnewline
\hline
Activity & 보이는 화면 하나 & 가게의 매장 앞 — 사용자가 볼 수 있고 만질 수 있어요 \tabularnewline
Service & 뒤에서 일하는 "보이지 않는 사람" & 창고 노동자 — 음악 재생, 파일 다운로드처럼 화면이 필요 없는 일을 해요 \tabularnewline
BroadcastReceiver & 시스템 브로드캐스트 수신 & 수위실 — "배터리가 다 떨어졌어요" "네트워크가 끊겼어요"를 알려 주는 확성기 \tabularnewline
ContentProvider & 앱 간 데이터 공유 & 공용 자료장 — 예를 들어 연락처를 다른 앱과 공유해요 \tabularnewline
\hline
\end{tabularx}
\end{center}

\begin{center}\small\color{gray}안드로이드 4대 컴포넌트 한눈에 보기\end{center}

초보 단계에서는 시간의 95\%를 Activity와 씨름하며 보내요. Service와 BroadcastReceiver는 음악 재생이나 메시지 푸시를 만들 때 쓰고, ContentProvider는 가장 적게 써요(자주 쓰는 데이터는 시스템이 이미 감싸 뒀으니까요). 한 가지 더 주의하세요. 4대 컴포넌트는 모두 AndroidManifest.xml에 “주민등록”을 해야 해요. 등록하지 않으면 시스템이 찾지 못해요.

\booksection{Activity 생명주기: 안드로이드 개발의 첫 번째 “보물지도”}

휴대폰의 앱은 언제든 끼어들 수 있어요. 전화가 오거나, 사용자가 위챗 답장하러 잠시 떠나거나, 시스템 메모리가 부족해 앱을 강제로 끝내 버리기도 하죠…… Activity의 일생은 바로 이 그림 위를 반복해 여행하는 거예요:

onCreate(탄생: 초기화)

onStart(보이기 시작)

onResume(포그라운드에서 상호작용 중)

onPause(일부 가려짐)

onStop(완전히 안 보임)

onDestroy(소멸)

onRestart / onStart / onResume: 포그라운드 복귀 시 되돌아가는 경로

\begin{center}\small\color{gray}그림 4: Activity 생명주기(진행 방향: 윗줄은 탄생 아랫줄은 퇴장, 돌아올 때는 왔던 길로 되돌아감)\end{center}

\begin{center}
\begin{tabularx}{\linewidth}{|X|X|X|}
\hline
\textbf{콜백} & \textbf{언제 호출되나} & \textbf{대표적인 용도} \tabularnewline
\hline
onCreate & Activity가 생성될 때(한 번) & 화면 로드, 데이터 초기화 \tabularnewline
onStart & 화면이 보이기 시작할 때 & 가벼운 로직 시작 \tabularnewline
onResume & 포커스를 얻고 상호작용 가능할 때 & 애니메이션 시작, 재생 재개 \tabularnewline
onPause & 일부 가려지거나 화면이 전환되는 순간 & 애니메이션 일시정지, 임시 저장 \tabularnewline
onStop & 완전히 보이지 않을 때 & 무거운 자원 해제, 다운로드 일시정지 \tabularnewline
onDestroy & 소멸되기 전 & 마지막 정리 \tabularnewline
\hline
\end{tabularx}
\end{center}

\begin{center}\small\color{gray}여섯 가지 생명주기 콜백의 대표적인 용도\end{center}

\textbf{주의: 화면 회전은 Activity를 “죽여요”}\par

초보자의 가장 고전적인 사고 현장: 앱 화면을 한 번 돌렸을 뿐인데 입력창의 글자가 모두 사라져요. 원인: 화면을 회전하면 기본적으로 Activity가 소멸 후 재생성되거든요(onDestroy onCreate를 한 바퀴 돌아요). 해법: rememberSaveable로 상태를 저장하세요(Compose에서 remember를 대체해요). 재생성될 때 값이 자동으로 복원돼요. 6장과 7장에서 계속 쓰게 될 거예요.

\booksection{Intent: 앱 안팎을 오가는 “차표”}

Intent(인텐트)는 컴포넌트끼리 소통하는 “차표”예요. 한 화면에서 다른 화면으로 넘어가거나, 브라우저를 열거나, 카메라를 띄우려면 모두 “표를 끊고 탑승”해야 해요:

\begin{lstlisting}[language=Kotlin]
// 방법 1: 명시적 Intent — 어떤 화면으로 갈지 직접 지정
val intent = Intent(this, DetailActivity::class.java)
intent.putExtra("name", 김철수)     // 차에 "짐"(데이터)을 싣는 셈
startActivity(intent)

// 방법 2: 암시적 Intent — 무엇을 할지만 말하면 시스템이 처리할 주체를 찾아요
val web = Intent(Intent.ACTION_VIEW, Uri.parse("https://developer.android.com"))
startActivity(web)   // 시스템이 브라우저를 열거나 선택하게 해요
\end{lstlisting}

두 방식은 쓰임이 달라요. 명시적(explicit)은 내 앱 안의 화면 전환에, 암시적(implicit)은 시스템이나 다른 앱의 기능을 호출할 때 써요(웹페이지 열기, 전화 걸기, 공유하기). Compose 프로젝트에서는 Composable 안에서 LocalContext.current로 Context를 얻으면 startActivity를 쓸 수 있어요.

\booksection{여러 화면과 내비게이션}

실제 앱에는 화면이 여러 개 있어요. Compose에서 공식 추천하는 내비게이션 방식은 Navigation Compose예요(의존성 라이브러리 androidx.navigation):

\begin{lstlisting}[language=Kotlin]
val navController = rememberNavController()

NavHost(navController, startDestination = "home") {
    composable("home") { HomeScreen(navController) }
    composable("detail/{name}") { entry ->
        val name = entry.arguments?.getString("name")
        DetailScreen(name ?: "알 수 없음")
    }
}

// 화면 안에서 이동:
navController.navigate("detail/김철수")     // 인자 전달
navController.popBackStack()              // 이전 화면으로 돌아가기
\end{lstlisting}

\begin{tcbinfo}{팁: 화면이 많아지면 어떻게 할까요?}
프로젝트가 커지면 모든 화면 라우트를 한 파일(예: AppNavGraph.kt)에 모아 통일해서 관리하세요. “어디로 이동할까”라는 로직이 여러 화면에 흩어지게 두지 말라는 뜻이에요. “돌아가기만 하는 코드”에서 “유지보수할 수 있는 코드”로 가는 첫걸음이에요.

\end{tcbinfo}

\booksection{곁가지로 알아두기: Fragment란 무엇일까(그리고 이 책에서 다루지 않은 이유)}

오래된 튜토리얼을 보다 보면 반드시 Fragment를 만나게 돼요. Activity 안에 들어가는 “재사용 가능한 화면 조각”으로, 구식 다중 화면 아키텍처의 절대 주인공이었죠. 하지만 Compose 시대에는 화면을 Composable 함수 + Navigation 라우트로 바로 만들면 충분하고, 공식적으로도 Fragment의 비중이 줄었어요. 그래서 이 책은 본줄기에 넣지 않았어요.

여러분을 위한 조언: 새 프로젝트에서는 Fragment를 굳이 쓰지 마세요. 다만 오래된 코드를 읽거나 면접에서 물어볼 때 “Activity에 들어가는 화면 조각”이라는 것만 알면 돼요. 학습 자료는 언제나 2023년 이후 것을 먼저 고르세요. 안드로이드 기술 스택은 생각보다 훨씬 빨리 낡아가니까요.

\begin{tcbpractice}{직접 해 보기: 카운터에 "기억력" 달기}
51장의 카운터를 개조해 봐요. (1) remember를 rememberSaveable로 바꾸고 화면을 회전해도 숫자가 남아 있는지 확인하세요. (2) “정보” 버튼을 추가해 Navigation으로 두 번째 화면으로 이동해 버전 번호를 보여 주세요. (3) 두 화면의 onCreate/onPause에 로그를 남기고(Log.d) Logcat으로 생명주기 호출 순서를 관찰하세요. 이 세 단계를 마치면 생명주기가 더 이상 종이 위의 이론이 아니에요.

\end{tcbpractice}

\begin{tcbsummary}{이번 장 정리}
\begin{itemize}
  \item 4대 컴포넌트: Activity(매장 얼굴), Service(뒷일 담당), BroadcastReceiver(확성기), ContentProvider(자료장)
  \item 생명주기 6대 콜백: onCreate onStart onResume onPause onStop onDestroy
  \item 화면 회전은 Activity를 소멸 후 재생성해요: rememberSaveable로 상태를 지키세요
  \item Intent는 "차표": 명시적은 내 앱 화면 이동, 암시적은 시스템 기능 호출
  \item 여러 화면은 Navigation Compose로: NavHost + composable 라우트
\end{itemize}

\end{tcbsummary}
